% !PRINTED 455
\weberpage{455}
%<<CONT
folgt, und wenn man wieder

\weq{29}{\beta c'-\beta c+\alpha p+\delta p=y}

setzt:

\weq{30}{\begin{aligned} 2(\beta c'+\alpha p) &= Bx+y,\\ 2(\beta c-\delta p) &= Bx-y. \end{aligned}}

Daraus

\weq{31}{y^2-Dx^2=4p^2.}

Aus (28) ergeben sich die Kongruenzen

\weq{32}{\begin{aligned} \beta c c' &\equiv Cx, \quad \beta(c+c') \equiv Bx, \quad \beta \equiv Ax\\ Dx^2 &\equiv \beta^2(c-c')^2 \pmod{p}. \end{aligned}}

Soll $c'$ von $c$ modulo $p$ verschieden sein, so kann hiernach weder $x$ noch $D$ durch $p$ teilbar sein und es folgt aus (32)

\weq{33}{\begin{aligned} A(c+c') &\equiv B,\\ Acc' &\equiv C, \end{aligned}}

und dies führt durch Elimination von $c'$ für $c$ auf die Kongruenz (13), also auf den auszuschließenden Wert von $c$.

\gesperrt{Fassen wir dies zusammen, so sehen wir, daß unter den Größen (15) nur in den beiden Ausnahmefällen $D=-3$, $D=-4$ äquivalente vorkommen, und zwar sind im Falle $D=-3$ je drei, im Falle $D=-4$ je zwei von ihnen äquivalent.}

In diesen Fällen muß also die rechte Seite der Formel (16) noch durch 3 oder durch 2 geteilt werden, und die Formel bleibt also auch dann noch richtig, wenn man, wie schon früher, für $D=-3, -4$ unter $h$ nicht 1, sondern $\frac{1}{3}$ und $\frac{1}{2}$ versteht.

\begin{center}
\textbf{\S\ 124. \gesperrt{Klassenkörper und Ordnungskörper.}}
\end{center}

Ist $D=Q^2\Delta$ eine Diskriminante und $\Delta$ ihr Stamm, so sind nach den Resultaten des vorigen Paragraphen die Klasseninvariante $u$ der Diskriminante $\Delta$ rational durch die Klasseninvariante $v$ der Diskriminante $D$ ausdrückbar, und zu jedem $u$ gehören $r$ Werte von $v$, wenn $r$ nach der Formel (16), \S\,123 bestimmt wird. Diese Werte von $v$ sind die Wurzeln einer Gleichung $r$ten Grades, deren Koeffizienten rational von $u$ abhängen. Die Größen $v$ gehören also einem Körper $\mathfrak{K}(D)$ über $\mathfrak{K}(\Delta)$ an, der kein anderer ist als der Klassenkörper der Diskriminante $D$.

In bezug auf die quadratischen Körper $\Omega=\mathfrak{R}(\sqrt{\Delta})$ sind beide Körper \gesperrt{Abelsche}. Wir wollen, wenn eine genauere Unter-
%CONT>>
